% ============================================================================
%  Capitolo 8 -- Elementi di algebra di Boole, parte 1           (Lezione 08)
% ============================================================================
\chapter{Algebra di Boole: operatori e proprietà}
\label{ch:boole1}

\epigraphline{Povero me! Povero me! Arriverò troppo tardi!}
  {Lewis Carroll, \emph{Alice nel Paese delle Meraviglie}, 1865}

\begin{chapgoals}
Alla fine di questo capitolo saprete
\begin{itemize}
\item descrivere un calcolatore come processore, memoria e ingresso/uscita, collegati da un bus;
\item spiegare la differenza tra l'analisi e la sintesi di un circuito logico;
\item scrivere le tabelle di verità di AND, OR, NOT, XOR, NAND e NOR e riconoscerne i simboli;
\item applicare le proprietà dell'algebra di Boole, leggi di De Morgan comprese, per semplificare un'espressione;
\item usare \ins{and}, \ins{or}, \ins{xor} e \ins{not} su \riscv per esaminare, porre a 1, azzerare e invertire singoli bit.
\end{itemize}
\end{chapgoals}
% Generated by tools/textbook/export_listings.js from tools/decks/mastery.js and tools/concepts/registry.json: do not edit.
\begin{chapstatus}{Stato di apprendimento di questo capitolo}
\masterychip{now}{DA COMPRENDERE OGGI}\enspace le maschere di bit; l'architettura di von Neumann; analisi e sintesi; l'algebra di Boole; le porte logiche.\par\smallskip
\masterychip{use}{DA UTILIZZARE OGGI}\enspace il programma completo, gli strumenti di PRG32 (si usa, non si studia, in questo corso).\par\smallskip
\masterychip{later}{DA COMPRENDERE PIÙ AVANTI}\enspace minimizzazione, Karnaugh (Capitolo~\ref{ch:boole2}).
\end{chapstatus}


% ----------------------------------------------------------------------------
\section{Riscaldamento sulla macchina}

State facendo algebra di Boole fin dai primi giorni. Il
Listato~\ref{lst:b1:left} è il modo in cui ogni gioco \prg chiede se una
direzione è premuta.

\asmfile[caption={\key{LEFT} è premuto?},label=lst:b1:left]{code/l08/c03.S}

La funzione \code{prg32\_input\_read} restituisce tutti i pulsanti insieme, un
bit per pulsante (Tabella~\ref{tab:b1:buttons}). Tenere premuti insieme
\key{RIGHT} e \key{A} dà $\hex{12} = 0001\,0010_2$. L'istruzione \ins{andi}
con \hex{01} conserva il bit~0 e forza a zero ogni altro bit, così che il
risultato è diverso da zero esattamente quando \key{LEFT} è premuto. La
costante usata con \ins{andi} si chiama \term{maschera}.

\begin{table}[htb]
\centering\small
\caption{Un bit per pulsante (le costanti \code{PRG32\_BTN\_*} di
\fname{prg32.h}).}\label{tab:b1:buttons}
\begin{tabular}{@{}lccccccc@{}}
\toprule
\textbf{Pulsante} & \key{LEFT} & \key{RIGHT} & \key{UP} & \key{DOWN} & \key{A} & \key{B} & \key{SELECT}\\
\midrule
Bit & 0 & 1 & 2 & 3 & 4 & 5 & 6\\
Maschera & \hex{01} & \hex{02} & \hex{04} & \hex{08} & \hex{10} & \hex{20} & \hex{40}\\
\bottomrule
\end{tabular}
\end{table}

\begin{figure}[htb]
\centering
\includegraphics[width=0.5\textwidth]{lab2_step1_lamps.png}
\caption{Sette lampade, una per bit della maschera dei pulsanti: una lampada è
accesa quando il suo bit vale 1.}\label{fig:b1:lamps}
\end{figure}

Ogni domanda sul joystick è un'espressione booleana su queste sette variabili.
E \ins{andi} è trentadue porte AND che lavorano fianco a fianco. Che cos'è,
esattamente, un AND?

% ----------------------------------------------------------------------------
\section{Il calcolatore come macchina logica}

Questa parte del libro scende di un livello sotto le istruzioni.

\subsection{Tre sottosistemi e un bus}

In termini odierni il modello di von Neumann del Capitolo~\ref{ch:history} si
disegna come nella Figura~\ref{fig:b1:bus}: le unità del 1945, collegate da
fili condivisi. Il \term{processore} (CPU) contiene l'unità di controllo, la
ALU e i registri, ed esegue le istruzioni. La \term{memoria centrale} contiene
le istruzioni e i dati. Il sottosistema di \term{ingresso/uscita} collega
joystick, display, suono, rete e memoria di massa. Il \term{bus} è l'insieme
dei fili che li unisce: linee degli indirizzi, linee dei dati e linee di
controllo. Tutto ciò che viaggia sul bus è una configurazione di bit.
Sull'ESP32-C6 i tre sottosistemi e il bus stanno tutti dentro un solo chip.

\begin{figure}[htb]
\centering
\begin{tikzpicture}
  \node[blk,fill=prgaccent,text=white,text width=36mm,minimum height=17mm] (cpu) at (0,0) {\textbf{Processore (CPU)}\\\footnotesize unità di controllo, ALU, registri};
  \node[blk,fill=prgsteel,text=white,text width=36mm,minimum height=17mm] (mem) at (4.7,0) {\textbf{Memoria centrale}\\\footnotesize istruzioni e dati};
  \node[blk,fill=prggreen,text=white,text width=36mm,minimum height=17mm] (io) at (9.4,0) {\textbf{Ingresso / uscita}\\\footnotesize joystick, display, suono, rete};
  \node[blk,fill=prgamber,text=white,minimum width=132mm,minimum height=7mm] (bus) at (4.7,-2.0) {\textbf{Bus}: linee degli indirizzi, dei dati, di controllo};
  \foreach \n in {cpu,mem,io} \draw[bus,<->] (\n.south) -- (\n.south|-bus.north);
\end{tikzpicture}
\caption{Il calcolatore elettronico: tre sottosistemi e un bus.}
\label{fig:b1:bus}
\end{figure}

\subsection{Il processore}

La capacità di elaborazione di un calcolatore risiede nel processore. Esso può
eseguire un insieme di azioni elementari, le \term{istruzioni}: comandi
espliciti che governano il trasferimento dell'informazione all'interno del
calcolatore e tra il calcolatore e i dispositivi di I/O, e che specificano le
operazioni aritmetiche e logiche da compiere. La memoria centrale contiene i
dati di ingresso e di uscita dell'elaborazione e la sequenza delle istruzioni;
una sequenza di istruzioni registrata in memoria centrale è un
\term{programma}. Il processore preleva le istruzioni dalla memoria e le esegue
una per una: \emph{prelievo, esecuzione, e da capo}.

Un calcolatore elettronico è, di conseguenza, un sistema

\begin{description}[style=nextline,leftmargin=1.5em]
\item[numerico:] un sistema digitale, in cui tutta l'informazione trattata è
  espressa come stringhe di bit;
\item[automatico:] una volta avviato procede da solo, istruzione dopo
  istruzione, senza intervento umano;
\item[a programma memorizzato:] il programma è registrato in memoria e può
  essere sostituito, così che la stessa macchina fa lavori diversi;
\item[elettronico:] è costruito con circuiti elettronici che memorizzano i bit
  e li elaborano.
\end{description}

\subsection{Circuiti logici}

I circuiti elettronici che memorizzano ed elaborano i bit si chiamano
\term{circuiti di commutazione} o \term{circuiti logici}. Un circuito logico
forma una \term{rete di commutazione}, e progettarne una si dice
\term{progetto logico}. Due livelli di tensione su un filo, due stati di un
interruttore, due valori di una variabile: per descrivere il comportamento di
questi circuiti si usa un'algebra, una notazione matematica che ne semplifica
sia l'analisi sia la sintesi. L'algebra è quella di \citet{boole1854},
applicata agli interruttori da Nakajima e da \citet{shannon1938}.

\begin{definition}{Analisi e sintesi}
\term{Analisi}: data la descrizione della \emph{struttura} di un circuito
(come è fatto), trovarne il \emph{comportamento} (che cosa fa).
\term{Sintesi}: data la descrizione del \emph{comportamento} (che cosa deve
fare), trovare una \emph{struttura} che lo realizzi (come costruirlo).
\end{definition}

Leggere uno schema e ricavarne la tabella di verità è analisi; lo è anche
leggere \code{andi t3, t0, 1} e dire che cosa finisce in \reg{t3}. Passare
dalla tabella di verità di una somma a un sommatore è sintesi; lo è anche
passare da «muoviti se \key{LEFT} o \key{RIGHT} è premuto» a una maschera e un
salto. La sintesi ha di solito molte soluzioni, e scegliere la meno costosa è
il problema della \term{minimizzazione}, argomento del
Capitolo~\ref{ch:boole2}.

% ----------------------------------------------------------------------------
\section{L'algebra di Boole}

\begin{definition}{Algebra di Boole}
L'\term{algebra di Boole} è l'algebra delle variabili che possono assumere solo
due valori, 0 e 1, con tre operazioni di base:
\begin{itemize}[nosep]
\item AND, il \term{prodotto logico} o congiunzione, scritto $x\cdot y$ o
  semplicemente $xy$;
\item OR, la \term{somma logica} o disgiunzione, scritta $x+y$;
\item NOT, il \term{complemento} o negazione, scritto $\nott{x}$ (nei
  programmi e sulla tastiera, \code{x'}).
\end{itemize}
\end{definition}

\subsection{Le tre operazioni fondamentali}

Un'operazione è definita dalla sua \term{tabella di verità}, che elenca
l'uscita per ogni combinazione degli ingressi (Figura~\ref{fig:b1:gates}). AND
vale 1 solo quando entrambi gli ingressi valgono 1; OR vale 0 solo quando
entrambi gli ingressi valgono 0; NOT inverte. Ogni operazione ha un simbolo
grafico, la \term{porta logica}, che sta per il circuito che la calcola.

\newcommand{\gatepanel}[4]{% port, title, formula, table body
\begin{minipage}[t]{0.31\textwidth}\centering
{\sffamily\bfseries\small\color{prgsteel}#2}\\[3pt]
\begin{circuitikz}\draw (0,0) node[#1](g){} (g.out) node[right,font=\small]{$z$}
  (g.in 1) node[left,font=\small]{$x$} (g.in 2) node[left,font=\small]{$y$};\end{circuitikz}\\[2pt]
{\small #3}\\[4pt]
{\small\ttfamily\begin{tabular}{cc|c}\textrm{$x$}&\textrm{$y$}&\textrm{$z$}\\\hline #4\end{tabular}}
\end{minipage}}

\begin{figure}[htb]
\centering
\gatepanel{and port}{AND (congiunzione)}{$z=x\cdot y$}{0&0&0\\0&1&0\\1&0&0\\1&1&1}\hfill
\gatepanel{or port}{OR (disgiunzione)}{$z=x+y$}{0&0&0\\0&1&1\\1&0&1\\1&1&1}\hfill
\begin{minipage}[t]{0.31\textwidth}\centering
{\sffamily\bfseries\small\color{prgsteel}NOT (negazione)}\\[3pt]
\begin{circuitikz}\draw (0,0) node[not port](g){} (g.out) node[right,font=\small]{$z$} (g.in) node[left,font=\small]{$x$};\end{circuitikz}\\[9pt]
{\small $z=\nott{x}$}\\[4pt]
{\small\ttfamily\begin{tabular}{c|c}\textrm{$x$}&\textrm{$z$}\\\hline 0&1\\1&0\end{tabular}}
\end{minipage}
\caption{Le tre operazioni fondamentali.}
\label{fig:b1:gates}
\end{figure}

\subsection{Altre tre porte}

Altre tre operazioni si incontrano dappertutto (Figura~\ref{fig:b1:gates2}).
XOR, l'\term{or esclusivo}, vale 1 quando gli ingressi sono diversi: è il bit
di somma di un'addizione, e \riscv lo ha come istruzione. NAND e NOR sono i
complementi di AND e OR; il cerchietto su un simbolo significa negazione.

\begin{figure}[htb]
\centering
\gatepanel{xor port}{XOR (or esclusivo)}{$z=\nott{x}y+x\nott{y}=x\oplus y$}{0&0&0\\0&1&1\\1&0&1\\1&1&0}\hfill
\gatepanel{nand port}{NAND (not and)}{$z=\nott{x\cdot y}$}{0&0&1\\0&1&1\\1&0&1\\1&1&0}\hfill
\gatepanel{nor port}{NOR (not or)}{$z=\nott{x+y}$}{0&0&1\\0&1&0\\1&0&0\\1&1&0}
\caption{XOR, NAND e NOR.}
\label{fig:b1:gates2}
\end{figure}

\subsection{Proprietà}

La Tabella~\ref{tab:b1:props} raccoglie le proprietà dell'algebra. Ogni
proprietà si presenta in coppia: una forma per AND e una per OR.

\begin{table}[htb]
\centering\small
\caption{Proprietà dell'algebra di Boole.}\label{tab:b1:props}
\begin{tabular}{@{}lll@{}}
\toprule
\textbf{Proprietà} & \textbf{Forma AND} & \textbf{Forma OR}\\
\midrule
Commutativa & $x\cdot y=y\cdot x$ & $x+y=y+x$\\
Associativa & $(x\cdot y)\cdot z=x\cdot(y\cdot z)$ & $(x+y)+z=x+(y+z)$\\
Idempotenza & $x\cdot x=x$ & $x+x=x$\\
Assorbimento & $x\cdot(x+y)=x$ & $x+(x\cdot y)=x$\\
Distributiva & $x\cdot(y+z)=x\cdot y+x\cdot z$ & $x+y\cdot z=(x+y)\cdot(x+z)$\\
Elemento neutro & $x\cdot 1=x$ & $x+0=x$\\
Elemento nullo & $x\cdot 0=0$ & $x+1=1$\\
Complemento & $x\cdot\nott{x}=0$ & $x+\nott{x}=1$\\
Involuzione & \multicolumn{2}{c}{$\nott{\nott{x}}=x$}\\
\bottomrule
\end{tabular}
\end{table}

Per la proprietà associativa, AND e OR si estendono a un numero qualsiasi di
operandi: $x\cdot y\cdot z$, $x+y+z$.

\begin{prgwarn}
Anche OR è distributivo rispetto ad AND: $x+yz=(x+y)(x+z)$. Questo è falso per
la somma e il prodotto ordinari, ed è la proprietà che si dimentica più spesso.
\end{prgwarn}

Come si dimostra una proprietà? Poiché una variabile ha solo due valori, ogni
identità può essere verificata provando tutti i casi: è il metodo
dell'\term{induzione perfetta}. Due espressioni sono \term{equivalenti} se
hanno la stessa tabella di verità. In alternativa una proprietà può essere
derivata algebricamente da altre; l'assorbimento, per esempio, segue da
elemento neutro, distributività ed elemento nullo:
\[
x+xy \;=\; x\cdot 1+x\cdot y \;=\; x\,(1+y) \;=\; x\cdot 1 \;=\; x .
\]

\subsection{Le leggi di De Morgan}

\begin{theorem}{De Morgan, 1847}
\[
\nott{x\cdot y}=\nott{x}+\nott{y},\qquad\qquad \nott{x+y}=\nott{x}\cdot\nott{y}.
\]
Il complemento di un prodotto è la somma dei complementi, e viceversa.
\end{theorem}

\noindent\emph{Dimostrazione} della prima legge, per induzione perfetta: le
colonne di $\nott{xy}$ e di $\nott{x}+\nott{y}$ sono identiche.

\begin{center}\small\ttfamily
\begin{tabular}{cc|c>{\columncolor{prgaccent!15}}c|cc>{\columncolor{prgaccent!15}}c}
\textrm{$x$}&\textrm{$y$}&\textrm{$xy$}&\textrm{$\nott{xy}$}&\textrm{$\nott{x}$}&\textrm{$\nott{y}$}&\textrm{$\nott{x}+\nott{y}$}\\\hline
0&0&0&1&1&1&1\\
0&1&0&1&1&0&1\\
1&0&0&1&0&1&1\\
1&1&1&0&0&0&0\\
\end{tabular}
\end{center}

La seconda legge si dimostra allo stesso modo (Esercizio~\ref{ex:b1:paper}).
\hfill$\square$

\smallskip
In termini di porte: un NAND è un OR degli ingressi negati, e un NOR è un AND
degli ingressi negati. Le leggi si generalizzano a un numero qualsiasi di
variabili.

\subsection{Dualità e precedenza}

\begin{theorem}{Principio di dualità}
Ogni identità dell'algebra di Boole resta vera se si scambiano AND con OR e 0
con 1.
\end{theorem}

Così $x\cdot 1=x$ diventa $x+0=x$, e ciascuna delle leggi di De Morgan è la
duale dell'altra. La dualità dimezza il lavoro: basta dimostrare una legge di
ogni coppia.

Quanto alla \term{precedenza}, NOT viene per primo, poi AND, poi OR, come per
potenze, prodotti e somme: $x+y\,\nott{z}$ significa
$x+(y\cdot(\nott{z}))$. Le parentesi cambiano l'ordine esattamente come
nell'algebra ordinaria.

\subsection{Semplificare con l'algebra}

Le proprietà sono strumenti per trasformare un'espressione in una equivalente
e meno costosa.

\begin{example}{Tre porte diventano un filo}
\begin{align*}
x\,y+x\,\nott{y} &= x\,(y+\nott{y}) && \text{proprietà distributiva, letta da destra a sinistra}\\
 &= x\cdot 1 && \text{legge del complemento}\\
 &= x && \text{elemento neutro.}
\end{align*}
Due porte AND, una porta OR e un invertitore sono diventati un pezzo di filo.
\end{example}

Allo stesso modo, usando la proprietà distributiva di OR rispetto ad AND,
\begin{equation}\label{eq:b1:simp}
x+\nott{x}\,y \;=\; (x+\nott{x})(x+y) \;=\; 1\cdot(x+y) \;=\; x+y ,
\end{equation}
un'identità chiamata talvolta \emph{teorema di semplificazione}, che sarà usata
di nuovo nel Capitolo~\ref{ch:boole2}.

% ----------------------------------------------------------------------------
\section{Ritorno alla macchina}

\subsection{Istruzioni bit a bit: l'algebra su un intero registro}

Ogni istruzione bit a bit di \riscv applica una porta a tutte le trentadue
coppie di bit insieme (Tabella~\ref{tab:b1:bitwise}, dove sono mostrati otto
bit). Le forme immediate \ins{andi}, \ins{ori} e \ins{xori} prendono una
costante di 12 bit. \ins{not} è una pseudoistruzione: \code{xori t0, a, -1},
perché $-1$ è tutti uni e $x\oplus 1=\nott{x}$.

\begin{table}[htb]
\centering\small
\caption{Istruzioni bit a bit, con $a=\texttt{1100\,1010}$ e
$b=\texttt{0000\,1111}$.}\label{tab:b1:bitwise}
\begin{tabular}{@{}llll@{}}
\toprule
\textbf{Istruzione} & \textbf{Operazione} & \textbf{Risultato} & \textbf{In C}\\
\midrule
\code{and t0, a, b} & AND bit a bit & \texttt{0000 1010} & \code{a \& b}\\
\code{or\ \ t0, a, b} & OR bit a bit & \texttt{1100 1111} & \code{a | b}\\
\code{xor t0, a, b} & XOR bit a bit & \texttt{1100 0101} & \code{a \textasciicircum\ b}\\
\code{not t0, a} & complementa ogni bit & \texttt{0011 0101} & \code{\textasciitilde a}\\
\bottomrule
\end{tabular}
\end{table}

\subsection{Quattro cose da fare con una maschera}

Le leggi dell'elemento neutro e dell'elemento nullo sono esattamente ciò che fa
funzionare una maschera (Tabella~\ref{tab:b1:mask}). Con AND, un 1 nella
maschera lascia passare il bit e uno 0 lo forza a 0; con OR, uno 0 lascia
passare il bit e un 1 lo forza a 1; con XOR, uno 0 lascia passare il bit e un
1 lo inverte.

\begin{table}[htb]
\centering\small
\caption{Quattro usi di una maschera (il bit 4 è il pulsante \key{A}).}\label{tab:b1:mask}
\begin{tabular}{@{}llll@{}}
\toprule
\textbf{Scopo} & \textbf{Legge usata} & \textbf{Istruzione} & \textbf{Altri bit}\\
\midrule
Esaminare un bit & $x\cdot 1=x$, $x\cdot 0=0$ & \code{andi t1, t0, 0x10} & forzati a 0 nel risultato\\
Porre un bit a 1 & $x+1=1$, $x+0=x$ & \code{ori\ \ t0, t0, 0x10} & invariati\\
Azzerare un bit & $x\cdot 0=0$, $x\cdot 1=x$ & \code{andi t0, t0, \textasciitilde 0x10} & invariati\\
Invertire un bit & $x\oplus 1=\nott{x}$, $x\oplus 0=x$ & \code{xori t0, t0, 0x10} & invariati\\
\bottomrule
\end{tabular}
\end{table}

La maschera \code{\textasciitilde 0x10} è \hex{FFFFFFEF}, che l'assemblatore
codifica come immediato $-17$.

\subsection{De Morgan al joystick}

Il verme è «fermo» quando né \key{LEFT} né \key{RIGHT} è premuto. Il
Listato~\ref{lst:b1:demorgan} scrive il test in due modi.

\asmfile[caption={$\nott{L+R}=\nott{L}\cdot\nott{R}$, come due programmi.},label=lst:b1:demorgan]{code/l08/c07.S}

Una maschera con due bit esamina due pulsanti in una sola istruzione: il
risultato è zero solo se entrambi i bit sono zero, cioè quando $\nott{L+R}$ è
vero. Per De Morgan questo è uguale a $\nott{L}\cdot\nott{R}$, che è la via
lunga scritta sotto. L'algebra dice che i due programmi sono equivalenti, e
quale costa meno: due istruzioni invece di cinque. La minimizzazione non
riguarda solo l'hardware.

\subsection{Una lampada per bit: scorrere e mascherare}

Le lampade della Figura~\ref{fig:b1:lamps} sono disegnate da un ciclo il cui
corpo è il Listato~\ref{lst:b1:lamp}. Per leggere il bit numero $n$ di una
parola, si fa scorrere la parola a destra di $n$ e si maschera con 1: è
l'idioma standard.

\asmfile[caption={Lettura del bit numero \reg{s0} della maschera dei pulsanti.},label=lst:b1:lamp]{code/l08/c08.S}

% ----------------------------------------------------------------------------
% Generated by tools/textbook/export_listings.js from decks/plan.js and code/l08.S: do not edit.
\section{Il programma di questo capitolo}\label{sec:lab:l08}

I listati di questo capitolo sono parti di un solo programma,
\fname{code/l08.S} nel repository del corso, stampato per intero nel
Listato~\ref{lst:lab:l08}. I test automatici del materiale del corso lo
compilano e lo eseguono sull'esecutore senza interfaccia di \prgqt{}, con
ingressi programmati, e controllano il frame; la
Figura~\ref{fig:lab:l08} è il frame scritto dal test.

\begin{description}[style=nextline,leftmargin=1.5em]
\item[Sullo schermo.] Sette lampade, una per pulsante (LEFT a sinistra), accese mentre il pulsante è premuto; un segnalino bianco mosso da LEFT e RIGHT; un quadrato verde quando né LEFT né RIGHT è premuto, rosso altrimenti.
\item[Che cosa premere.] LEFT e RIGHT muovono il segnalino; provate anche A e B.
\end{description}

\begin{figure}[htb]
\centering
\includegraphics[width=0.5\textwidth]{lab_l08.png}
\caption{\fname{l08.S} su \prgqt{}.}\label{fig:lab:l08}
\end{figure}

\asmfile[caption={Il programma completo del capitolo, \fname{code/l08.S}.},label=lst:lab:l08]{code/l08/program.S}

Per eseguirlo su \prgqt{} i comandi sono quelli del
Listato~\ref{lst:lab:l08:run}; il Paragrafo~\ref{sec:prg:testing} dà gli
stessi passi per QEMU e per la scheda.

\begin{shell}[caption={Compilare \fname{l08.S} ed eseguirlo su \prgqt{}.},label=lst:lab:l08:run]
# 1. start PRG32-QT on this computer (it listens on port 8080)

# 2. in the PRG32 folder, with ESP-IDF active: build the cartridge
python3 -m prg32 cartridge build ~/arc12/code/l08.S --portable \
    --entry-prefix l08 --name l08 --out l08.prg32

# 3. send it to PRG32-QT and start it
python3 -m prg32 esp32c6 upload-and-run l08.prg32 \
    --url http://localhost:8080 --slot cart0

# 4. look, press, compare with the previous slide
\end{shell}


% ----------------------------------------------------------------------------
\section{Verificate di aver capito}

\begin{quiz}{Per quali ingressi $x$ NAND $y$ vale 0?}
\item $x=0$, $y=0$
\item $x=0$, $y=1$
\item $x=1$, $y=1$
\item mai
\end{quiz}

\begin{quiz}{Per De Morgan, $\nott{a\cdot\nott{b}}$ è uguale a:}
\item $\nott{a}\cdot b$
\item $\nott{a}+b$
\item $a+\nott{b}$
\item $\nott{a}+\nott{b}$
\end{quiz}

\begin{quiz}{$\reg{t0}=\hex{0000005A}$. Che cosa c'è in \reg{t1} dopo \code{andi t1, t0, 0x0F}?}
\item \hex{0000005A}
\item \hex{00000050}
\item \hex{0000000A}
\item \hex{0000000F}
\end{quiz}

\begin{quiz}{Quale istruzione porta a 1 il bit 5 di \reg{t0} e lascia invariati tutti gli altri bit?}
\item \code{andi t0, t0, 0x20}
\item \code{ori t0, t0, 0x20}
\item \code{xori t0, t0, 0x20}
\item \code{addi t0, t0, 0x20}
\end{quiz}

\section{Esercizi}

\begin{exercise}[prima su carta, poi su \prg]\label{ex:b1:paper}
\begin{enumerate}[nosep]
\item Dimostrate la seconda legge di De Morgan,
  $\nott{x+y}=\nott{x}\cdot\nott{y}$, con una tabella di verità.
\item Semplificate con l'algebra $x\,y+x\,\nott{y}\,z+\nott{x}\,y$, nominando
  la legge usata a ogni passo.
\item Scrivete la maschera e le istruzioni che verificano se \key{A} e
  \key{B} sono premuti \emph{insieme}.
\item Modificate il programma delle lampade in modo che le lampade mostrino il
  complemento dei pulsanti (accese quando rilasciati).
\item Convertite una lettera da maiuscola a minuscola e viceversa usando una
  sola istruzione bit a bit (suggerimento: la tabella ASCII del
  Capitolo~\ref{ch:information}).
\end{enumerate}
\expected{Su carta, il secondo compito dà
$y\,(x+\nott{x})+x\,\nott{y}\,z = y+x\,\nott{y}\,z = y+x\,z$, l'ultimo passo
per la (\ref{eq:b1:simp}). Per il terzo: \code{andi t1, t0, 0x30} e poi un
confronto \emph{con} \hex{30}; verificare che sia diverso da zero
significherebbe \key{A} \emph{oppure} \key{B}. Su \prg: tutte e sette le
lampade accese a riposo, ciascuna che si spegne mentre il suo pulsante è tenuto
premuto. Per la lettera, \ins{xori} con \hex{20} inverte il bit~5.}
\end{exercise}

\section{Glossario}

\begin{glossarytable}
logic circuit & circuito logico\\
switching network & rete di commutazione\\
analysis / synthesis & analisi / sintesi\\
logical product (AND) & prodotto logico (congiunzione)\\
logical sum (OR) & somma logica (disgiunzione)\\
involution & involuzione\\
truth table & tabella di verità\\
logic gate & porta logica\\
complement & complemento (negazione)\\
duality & dualità\\
mask & maschera\\
to set / clear / toggle & porre a 1 / azzerare / invertire\\
\end{glossarytable}

\section{Conclusioni e letture}

L'architettura di von Neumann rappresenta una macchina reale in forma
semplificata: processore, memoria, ingresso/uscita. L'algebra di Boole
rappresenta funzioni logiche che si possono realizzare con circuiti
elettronici. Tre operazioni, una manciata di proprietà e le leggi di De Morgan
bastano a trasformare qualsiasi espressione. Un'istruzione bit a bit è
l'algebra applicata a trentadue bit insieme, e una maschera è la legge
dell'elemento neutro al lavoro.

\paragraph{Letture.} \citet{patterson2020}, appendice~A, «The Basics of Logic
Design»; \emph{Playing with RISC-V}, capitolo~6, «State and Input: Branches,
Bitmasks, and Boundaries» \citep{prg32textbook}; \citet{bucci2017}, i capitoli
sull'algebra di Boole.
